% ══════════════════════════════════════════════════════════════════════════════
% exercise_subgroup_analysis.tex
% Descriptive subgroup heterogeneity for wearable-defined exercise episodes.
% All figures generated by scripts/render_exercise_thesis_figure_suite.py
% and saved to outputs/exercise_thesis_figures_2/.
% ══════════════════════════════════════════════════════════════════════════════

\section{Descriptive Heterogeneity of the Exercise Glucose Response}
\label{sec:exercise_subgroup}

This section characterises descriptive heterogeneity in the glucose response
to wearable-defined ambulatory exercise-like episodes across participant
subgroups, cadence--strain combinations, and metabolic phenotypes.
All analyses are associational; no subgroup contrast is a causal effect.
Groups with fewer than 20 episodes or 10 unique participants are suppressed.
Legend entries report the number of episodes~(\emph{n}) and unique
participants~(\emph{p}); for quartile-based subgroups, the effective
within-group numeric range is also shown.

% ── Overall kernel ────────────────────────────────────────────────────────────

\subsection{Overall Glucose-Response Kernel}
\label{subsec:ex_kernel_overall}

\begin{figure}[H]
  \centering
  \includegraphics[width=0.97\linewidth]{exercise_v2/exercise_glucose_response_overall.pdf}
  \caption{%
    \textbf{Overall wearable-defined exercise glucose kernel.}
    Descriptive glucose kernel around wearable-defined ambulatory
    exercise-like episodes, with participant-clustered bootstrap 95\,\%
    intervals. Curves are additionally stratified by moderate and vigorous
    strain where eligible.
    This is not a causal exercise estimate.%
  }
  \label{fig:ex_kernel_overall}
\end{figure}

% ── Cadence × strain heatmap ─────────────────────────────────────────────────

\subsection{Cadence and Strain Heterogeneity}
\label{subsec:ex_cadence_strain}

\begin{figure}[H]
  \centering
  \includegraphics[width=0.97\linewidth]{exercise_v2/exercise_response_by_cadence_strain.pdf}
  \caption{%
    \textbf{Exercise glucose response by cadence and strain class.}
    Left: episode (\emph{n}) and participant (\emph{p}) counts.
    Centre: observed glucose change from refined onset to $+60$\,min
    ($\Delta G_{60}$, mg/dL).
    Right: observed minus frozen-model prediction at $+60$\,min (signed
    residual, mg/dL).
    Cells below the eligibility threshold are marked ``insufficient n''.
    Run-like cadence is folded into brisk walk when sparse.
    Text colour adapts to cell luminance for readability.%
  }
  \label{fig:ex_cadence_strain}
\end{figure}

% ── Glycemic stratum ──────────────────────────────────────────────────────────

\subsection{Glycemic Stratum}
\label{subsec:ex_glycemic}

\begin{figure}[H]
  \centering
  \includegraphics[width=0.97\linewidth]{exercise_v2/exercise_response_by_glycemic_stratum.pdf}
  \caption{%
    \textbf{Exercise glucose response by glycemic stratum.}
    Descriptive heterogeneity across normoglycaemia, prediabetes,
    T2D oral/non-insulin, and T2D insulin-dependent strata.
    Right panel: frozen-model residual trajectory (observed $-$ prediction)
    from the matched residual diagnostic; available only for non-insulin strata.
    Differences are not causal effects of diabetes status.%
  }
  \label{fig:ex_glycemic}
\end{figure}

% ── Baseline glucose ──────────────────────────────────────────────────────────

\subsection{Baseline Glucose Dependence}
\label{subsec:ex_baseline_glucose}

\begin{figure}[H]
  \centering
  \includegraphics[width=0.97\linewidth]{exercise_v2/exercise_response_by_baseline_glucose.pdf}
  \caption{%
    \textbf{Exercise glucose response by baseline glucose at episode onset.}
    Left: $\Delta G$ trajectory by 30\,mg/dL baseline bin.
    Centre: observed glucose level after onset.
    Right: LOWESS trend of within-episode glucose nadir vs.\ baseline glucose.
    Baseline dependence is descriptive and may reflect regression to the mean,
    safety gating, and selection, not a causal dose-response.%
  }
  \label{fig:ex_baseline_glucose}
\end{figure}

% ── Continuous quartile subgroups ─────────────────────────────────────────────

\subsection{Anthropometric and Metabolic Quartile Subgroups}
\label{subsec:ex_quartile_subgroups}

Each figure below uses a three-panel layout:
left panel shows the $\Delta G$ response kernel over time (0--180\,min
after refined onset) with participant-clustered bootstrap 95\,\% bands;
centre panel shows $\Delta G_{60}$ (filled circles) and minimum $\Delta G$
(open squares) with 95\,\% bootstrap error bars;
right panel shows the observed minus frozen prediction at $+60$\,min.
Legend labels include the quartile name, effective participant-level value
range (with unit), episode count~(\emph{n}), and participant count~(\emph{p}).

\begin{figure}[H]
  \centering
  \includegraphics[width=0.97\linewidth]{exercise_v2/exercise_response_by_age_quartile.pdf}
  \caption{%
    \textbf{Age quartile: exercise glucose response kernel.}
    Descriptive subgroup comparison across age quartiles (unit: years).
    Only quartiles with $\geq 20$ episodes and $\geq 10$ participants are shown.%
  }
  \label{fig:ex_age_quartile}
\end{figure}

\begin{figure}[H]
  \centering
  \includegraphics[width=0.97\linewidth]{exercise_v2/exercise_response_by_bmi_quartile.pdf}
  \caption{%
    \textbf{BMI quartile: exercise glucose response kernel.}
    Descriptive subgroup comparison across BMI quartiles (unit: kg/m$^2$).%
  }
  \label{fig:ex_bmi_quartile}
\end{figure}

\begin{figure}[H]
  \centering
  \includegraphics[width=0.97\linewidth]{exercise_v2/exercise_response_by_hba1c_quartile.pdf}
  \caption{%
    \textbf{HbA1c quartile: exercise glucose response kernel.}
    Descriptive subgroup comparison across HbA1c quartiles (unit: \%).%
  }
  \label{fig:ex_hba1c_quartile}
\end{figure}

\begin{figure}[H]
  \centering
  \includegraphics[width=0.97\linewidth]{exercise_v2/exercise_response_by_c_peptide_quartile.pdf}
  \caption{%
    \textbf{C-peptide quartile: exercise glucose response kernel.}
    Descriptive subgroup comparison across C-peptide quartiles
    (unit: ng/mL).%
  }
  \label{fig:ex_c_peptide_quartile}
\end{figure}

\begin{figure}[H]
  \centering
  \includegraphics[width=0.97\linewidth]{exercise_v2/exercise_response_by_resting_hr_quartile.pdf}
  \caption{%
    \textbf{Resting heart-rate quartile: exercise glucose response kernel.}
    Descriptive subgroup comparison across resting HR quartiles (unit: bpm).%
  }
  \label{fig:ex_resting_hr_quartile}
\end{figure}

\begin{figure}[H]
  \centering
  \includegraphics[width=0.97\linewidth]{exercise_v2/exercise_response_by_waist_to_hip_ratio_quartile.pdf}
  \caption{%
    \textbf{Waist-to-hip ratio quartile: exercise glucose response kernel.}
    Descriptive subgroup comparison across waist-to-hip ratio quartiles
    (dimensionless).%
  }
  \label{fig:ex_whr_quartile}
\end{figure}

\begin{figure}[H]
  \centering
  \includegraphics[width=0.97\linewidth]{exercise_v2/exercise_response_by_systolic_bp_quartile.pdf}
  \caption{%
    \textbf{Systolic blood-pressure quartile: exercise glucose response kernel.}
    Descriptive subgroup comparison across systolic BP quartiles (unit: mmHg).%
  }
  \label{fig:ex_systolic_bp_quartile}
\end{figure}

\begin{figure}[H]
  \centering
  \includegraphics[width=0.97\linewidth]{exercise_v2/exercise_response_by_diastolic_bp_quartile.pdf}
  \caption{%
    \textbf{Diastolic blood-pressure quartile: exercise glucose response kernel.}
    Descriptive subgroup comparison across diastolic BP quartiles (unit: mmHg).%
  }
  \label{fig:ex_diastolic_bp_quartile}
\end{figure}

% ── Categorical subgroups ─────────────────────────────────────────────────────

\subsection{Categorical Subgroups}
\label{subsec:ex_categorical_subgroups}

\begin{figure}[H]
  \centering
  \includegraphics[width=0.97\linewidth]{exercise_v2/exercise_response_by_sex_at_birth.pdf}
  \caption{%
    \textbf{Sex at birth: exercise glucose response kernel.}%
  }
  \label{fig:ex_sex_at_birth}
\end{figure}

\begin{figure}[H]
  \centering
  \includegraphics[width=0.97\linewidth]{exercise_v2/exercise_response_by_ethnicity.pdf}
  \caption{%
    \textbf{Ethnicity: exercise glucose response kernel.}
    Categories derived from the four AI-READI ethnicity indicator columns
    (\texttt{hispanic\_or\_latino}, \texttt{not\_hispanic\_or\_latino},
    \texttt{dont\_know}, \texttt{prefer\_not\_to\_answer}).
    Groups with insufficient episodes are suppressed.%
  }
  \label{fig:ex_ethnicity}
\end{figure}

\begin{figure}[H]
  \centering
  \includegraphics[width=0.97\linewidth]{exercise_v2/exercise_response_by_race.pdf}
  \caption{%
    \textbf{Race: exercise glucose response kernel.}
    Derived from the primary race field; ``Prefer not to say'' and
    ``Not answered'' are unified as ``Unknown / not reported''.%
  }
  \label{fig:ex_race}
\end{figure}

\begin{figure}[H]
  \centering
  \includegraphics[width=0.97\linewidth]{exercise_v2/exercise_response_by_clinical_site.pdf}
  \caption{%
    \textbf{Clinical site: exercise glucose response kernel.}%
  }
  \label{fig:ex_clinical_site}
\end{figure}

\begin{figure}[H]
  \centering
  \includegraphics[width=0.97\linewidth]{exercise_v2/exercise_response_by_med_insulin.pdf}
  \caption{%
    \textbf{Insulin medication: exercise glucose response kernel.}
    Episodes stratified by whether the participant was on insulin therapy
    at the time of the study.%
  }
  \label{fig:ex_med_insulin}
\end{figure}

\begin{figure}[H]
  \centering
  \includegraphics[width=0.97\linewidth]{exercise_v2/exercise_response_by_med_any_diabetes_drug.pdf}
  \caption{%
    \textbf{Any diabetes medication: exercise glucose response kernel.}
    Episodes stratified by whether the participant was on any
    glucose-lowering medication.%
  }
  \label{fig:ex_med_any_diabetes_drug}
\end{figure}

% ── Residual figures ──────────────────────────────────────────────────────────

\subsection{Frozen-Model Residual Summary}
\label{subsec:ex_residual_summary}

\begin{figure}[H]
  \centering
  \includegraphics[width=0.97\linewidth]{exercise_v2/exercise_residual_response_matched_controls.pdf}
  \caption{%
    \textbf{Frozen-model residual trajectory: episodes vs.\ matched controls.}
    Signed residual is observed glucose minus frozen-model prediction.
    Positive values indicate the frozen model underpredicts observed glucose.
    Confidence intervals are participant-clustered bootstrap 95\,\% bands.%
  }
  \label{fig:ex_residual_matched}
\end{figure}

\begin{figure}[H]
  \centering
  \includegraphics[width=0.97\linewidth]{exercise_v2/exercise_mae_bias_summary_by_strain.pdf}
  \caption{%
    \textbf{Forecast error summary by strain class.}
    Left: mean absolute error (MAE, mg/dL) for exercise-like episodes and
    matched controls.
    Right: signed residual (observed $-$ prediction, mg/dL).
    Light-strain estimates are sparse and are not interpreted.%
  }
  \label{fig:ex_mae_bias}
\end{figure}

% ── Planning-head diagnostics ─────────────────────────────────────────────────

\subsection{Exercise Planning Head: Response Surface and Horizon Shapes}
\label{subsec:ex_head_diagnostics}

\begin{figure}[H]
  \centering
  \includegraphics[width=0.80\linewidth]{exercise_v2/exercise_head_response_surface.pdf}
  \caption{%
    \textbf{Imposed planning-head response surface.}
    The surface shows $\Delta\hat{y}_{\mathrm{ex}} = -A\,\mathrm{ReLU}(\hat{y}_{\mathrm{base}}-105)\,S$
    as a function of base-forecast glucose and HR activation~$A$.
    The safety floor at 105\,mg/dL suppresses the overlay at low glucose.
    Magnitude is imposed, not identified from AI-READI data.%
  }
  \label{fig:ex_head_surface}
\end{figure}

\begin{figure}[H]
  \centering
  \includegraphics[width=0.80\linewidth]{exercise_v2/exercise_head_horizon_shapes.pdf}
  \caption{%
    \textbf{Imposed planning-head horizon shapes.}
    Three canonical scenarios: low glucose with moderate HR (safety gate active),
    high glucose with moderate HR, and high glucose with vigorous HR.
    No causal claim; the shape is prescribed by the planning head formula.%
  }
  \label{fig:ex_head_horizons}
\end{figure}
